본문은 그냥 쓰시면 됩니다.

\subsection{이미지 첨부 방법}
\begin{figure}[h]
    \centering
    \includegraphics[width=0.75\linewidth]{figure/c_lanigera.jpg}
    \caption{성체 Chinchilla lanigera}
    \label{fig:firstImage}
\end{figure}
이미지는 이렇게 Figure \ref{fig:firstImage} 가져다가 쓸 수 있다.
이미지 번호는 그냥 알아서 해준다.

\subsection{표 첨부 방법}
\begin{table}[ht]
    \centering
    \caption{Booktabs를 사용한 표}
    \label{tab:booktabs}
    \begin{tabular}{lcc} % 세로줄 없음. l: 왼쪽 정렬, c: 중앙 정렬
        \toprule % 가장 위쪽 굵은 선
        \textbf{항목} & \textbf{값} & \textbf{단위} \\
        \midrule % 중간 선
        길이 & $10.5$ & $\text{cm}$ \\
        질량 & $200.0$ & $\text{g}$ \\
        \bottomrule % 가장 아래쪽 굵은 선
    \end{tabular}
\end{table}


\subsection{수도코드 첨부 방법}
수도코드는 아래처럼 하면 된다.

\begin{algorithm}[!ht]
    \caption{Your Algorithm Name}
    \label{alg:your_algorithm}
    \begin{algorithmic}[1] % The number tells where the line numbering should start
        \REQUIRE $n \geq 0 \vee x \neq 0$ % Input requirements
        \ENSURE $y = x^n$ % Output specifications
        \STATE $y \leftarrow 1$
        \STATE $X \leftarrow x$
        \STATE $N \leftarrow n$
        \WHILE{$N \neq 0$}
            \IF{$N$ is even}
                \STATE $X \leftarrow X \times X$
                \STATE $N \leftarrow N / 2$ 
            \ELSE[$N$ is odd]
                \STATE $y \leftarrow y \times X$
                \STATE $N \leftarrow N - 1$
            \ENDIF
        \ENDWHILE
    \end{algorithmic}
\end{algorithm}